\documentclass[12pt]{article}
\usepackage[utf8]{inputenc}
\usepackage[T1]{fontenc}
\usepackage{hyperref}
\title{Phupha Language Resources for K-12 Curriculum Development}
\author{Compiled by Claude Sonnet 5.5 (claude-sonnet-5-5) for Andr\'as Kornai}
\date{}
\begin{document}
\maketitle

\section{Introduction}
Phupha (ISO 639-3 yph, Glottocode phup1237; also called Downriver Phula) is a Loloish language of Yunnan, China. Wikipedia describes four closely related varieties that may be distinct languages: Phupha, Alugu (aub), Phupa (ypp) and Phuza (ypz). Phupha itself is spoken in southwest Gejiu County (4 villages) and one village in Yuanyang County. Population figures differ by source. Wikipedia gives 14,000 (2007) for the whole cluster. Global Recordings Network lists 1,600 for Phupha, and Joshua Project gives 1,600 for yph. The brief gives about 1,300. Wikipedia reports that use is declining, except for Alugu, which it says is taught in primary schools. That is a different variety and does not mean Phupha is taught. Joshua Project notes growing use of Mandarin, especially among the young.

\section{Primary Educational Resources}
No primers, textbooks or school programmes for Phupha turned up in the searches. The language is not known to be taught in schools.

\section{Linguistic and Reference Materials}
\begin{itemize}
\item Phupha language (Wikipedia): \url{https://en.wikipedia.org/wiki/Phupha_language}
\end{itemize}
Other Phula sources were not found in the Phupha-specific search.

\section{Religious and Community Resources}
\begin{itemize}
\item Joshua Project profile of Phupha in China: \url{https://joshuaproject.net/people_groups/20628}
\item Joshua Project profile of Phupa in China (a related variety, ypp): \url{https://joshuaproject.net/people_groups/20611}
\end{itemize}

\section{Digital and Audio Resources}
\begin{itemize}
\item Global Recordings Network language page for yph: \url{https://globalrecordings.net/en/language/yph}
\end{itemize}
I did not verify what recordings the GRN page holds.

\section{Curriculum Development Context}
Searches found no orthography, dictionary, grammar or classroom materials for Phupha. With about 1,300 to 1,600 speakers and rising Mandarin use among the young, documentation and intergenerational transmission come before any curriculum. The Alugu primary-school teaching reported by Wikipedia may be worth investigating as a regional precedent. I did not find details of it.
\end{document}